\documentclass[10pt,a4paper]{amsart}
\usepackage[margin=19mm]{geometry}
\usepackage{amsmath,amssymb,mathtools}
\usepackage[scr=rsfs,cal=euler]{mathalfa}
\usepackage{xfrac}
\usepackage[unicode,hidelinks,bookmarks=false]{hyperref}
\newcommand{\ie}{i.e.\ }
\newcommand{\cf}{cf.\ }
\newcommand{\resp}{respectively\ }
\newcommand{\Subsection}{\subsection}
\providecommand{\nobreakdash}{\nobreak\hbox{-}}
\setlength{\emergencystretch}{2em}
\multlinegap0pt
\usepackage{amssymb}
\usepackage{mathtools}
\usepackage{csquotes}
\usepackage{tikz}
\newtheorem{lem}{Lemma}
\newcommand{\bu}{_{\bullet}}
\title{Persistent Cost of Lipschitz Maps}
\author{Francisco J. Gozzi, Manuela A. Cerdeiro, Pablo E. Riera}
\date{Source: arXiv:2511.07674v2 (CC BY 4.0)}
\begin{document}
\maketitle

\noindent\emph{Source and licence.} Persistent Cost of Lipschitz Maps. Version arXiv:2511.07674v2 (CC BY 4.0), distributed by arXiv under the Creative Commons Attribution 4.0 International licence, http://creativecommons.org/licenses/by/4.0/, as recorded in the arXiv licence metadata for this exact version and shown on the abstract page https://arxiv.org/abs/2511.07674v2. Full licence notice: \texttt{licenses/arxiv-2511.07674-CC-BY-4.0.txt}.\par\smallskip

\noindent\textit{Editorial scope.} Counted unit: the article's lem homology (source0.tex lines 442--448). The article's complete original proof of it is delivered with it (source0.tex lines 450--472, one \texttt{proof} environment of 23 lines). No mathematics is written, repaired or re-derived: the statement and the proof are the article's own lines, in the article's own notation, and no retained line is substituted or re-flowed.  No further source-local context is needed: every cross-reference of every retained line resolves inside the two delivered spans.  Nothing was omitted from any retained span, so the package carries no omission record.  No retained line contains a citation, so the wrapper carries no bibliography and no bibliography entry.  No retained line contains a figure environment, an image-inclusion command or drawing code, so no image is delivered and the package declares no asset dependency.  Editorial formatting only, in addition: an AMS \texttt{amsart} preamble that restates the article's own declarations and macros used by the retained lines, namely the environments \texttt{lem} on separate counters, together with the standard packages those lines need.  The retained environments print in this document as Lemma 1 in order of appearance, whereas the article prints the same items with the numbering of its own full document.  The complete native source of the article as distributed in the arXiv e-print for this exact version is bundled unchanged as \texttt{sources/188827/source0.tex}.
\begin{lem}\label{lemma induced function on homology}
Let \( h : X \to Y \) be an \( \epsilon \)-quasi \( 1 \)-Lipschitz map. Then \( h \) induces an \( \epsilon \)-degree persistence homomorphism: $
h_* : H_d(\operatorname{VR}\bu(X)) \to H_d(\operatorname{VR}_{\bullet+\epsilon}(Y)).$


Moreover, if two \( \epsilon \)-quasi \( 1 \)-Lipschitz maps \( h_1, h_2 : X \to Y \) are within uniform distance \( \delta \geq 0 \), then they induce the same \( (\epsilon + \delta) \)-degree map on homology.
\end{lem}

\begin{proof}
The map \( h \) induces a simplicial map:
\[
\hat{h} : \operatorname{VR}_{\bullet}(X) \to \operatorname{VR}_{\bullet+\epsilon}(Y), \quad [x_0 \dots x_r] \mapsto [h(x_0) \dots h(x_r)],
\]
since a set \( \sigma \subseteq X \) of diameter \( \ell \) maps to a set \( h(\sigma) \subseteq Y \) of diameter at most \( \ell + \epsilon \). This commutes with inclusions and thus gives a homomorphism between filtrations. The homology functor yields the desired persistence homomorphism.

For the second claim, suppose \( d_Y(h_1(x), h_2(x)) \leq \delta \) for all \( x \in X \). For any simplex \( \sigma \in \operatorname{VR}_{\bullet}(X) \) of diameter \( \ell \), consider the union:
\[
h_1(\sigma) \cup h_2(\sigma) \subseteq Y.
\]
For any \( y, y' \in h_1(\sigma) \cup h_2(\sigma) \), we have:
- If both come from the same \( h_i \), then \( d(y, y') \leq \ell + \epsilon \).
- If \( y = h_1(x) \), \( y' = h_2(x') \), then:
\[
d(y, y') \leq d(h_1(x), h_2(x)) + d(h_2(x), h_2(x')) \leq \delta + (\ell + \epsilon).
\]
Thus, the diameter of \( h_1(\sigma) \cup h_2(\sigma) \) is at most \( \ell + \epsilon + \delta \). This implies that \( \hat{h}_1 \) and \( \hat{h}_2 \) are contiguous as maps:
\[
\operatorname{VR}_{\bullet}(X) \to \operatorname{VR}_{\bullet+\epsilon+\delta}(Y),
\]
and therefore induce the same map on homology.
\end{proof}

\end{document}
